\section{GlyphNet Unified Inbox: Encrypted Pilot-to-Pilot Work}

The first Unified Inbox vertical slice connects Pilot's existing personal task authority to an
encrypted mother-to-mother transport.  It proves the intended product behaviour: one person can
ask Pilot to give another person a task; the recipient's Pilot can receive it privately; and the
recipient decides whether it becomes their task.  No browser development token, static Wave
identity or shared television utterance grants messaging authority.

\subsection{Authority boundaries}

The implementation deliberately preserves two separate responsibilities.  The established
Pilot Inbox remains authoritative for personal lists, tasks, reminders and recipient
acceptance.  The Unified Inbox bridge is responsible for encrypted transport, message
envelopes, delivery receipts, replay prevention and reconciliation between mother brains.  A
verified incoming proposal is imported into the recipient's Pilot Inbox exactly once and remains
awaiting recipient acceptance until that recipient acts from a trusted phone.

Every private operation is bound to the mobile trust system described in the preceding section.
The phone certificate and current capability lease must validate, the requested scope must be
present, the request persona must match the certificate persona, and the phone must sign the exact
canonical request.  Removing the phone from Production Private Identity therefore blocks later
message, task and response operations immediately.

\subsection{Encrypted peer relationship}

Each mother brain maintains two independent key roles:
\begin{itemize}
  \item its existing Ed25519 identity signs descriptors, packets and receipts; and
  \item a dedicated X25519 transport key establishes encryption keys for Inbox payloads.
\end{itemize}

A peer descriptor contains the mother identifier, signing public key, transport public key and
issue time.  The descriptor is Ed25519 signed and verified before it enters the local trusted-peer
registry.  Private keys remain in owner-only local files.

Each outbound packet uses a new ephemeral X25519 key.  The sender combines that key with the
recipient mother's transport public key, then derives a 256-bit content key through HKDF-SHA-256.
The message or structured task is encrypted with AES-256-GCM.  Packet routing metadata is supplied
as authenticated additional data, and the complete packet is signed by the sending mother.  The
recipient verifies the packet hash, signature, intended mother and intended persona before
decryption.  Changing any destination, ciphertext or routing field causes the operation to fail.

The local Unified Inbox state also encrypts message bodies and task titles with AES-256-GCM.
Receipts contain identifiers, states, timestamps and evidence hashes, never message text or task
content.  The existing personal task store still retains the accepted task information required
to operate the user's task list; it is not copied into a public chain or transport journal.

\subsection{Typed-message flow}

For an ordinary Pilot-to-Pilot message, the sender's phone signs an idempotent request containing
the exact sender, recipient, recipient mother and message body.  The sender mother checks
\texttt{message.send}, encrypts one recipient-specific packet and records the local encrypted
envelope.  The receiving mother verifies, decrypts and stores one encrypted local copy.  Repeating
the same request returns the original message; reusing its idempotency key with different content
is rejected.

\subsection{Structured task flow}

The task journey is intentionally more explicit:
\begin{enumerate}
  \item The sender chooses one resolved Pilot recipient and describes the task.
  \item Pilot creates an encrypted prepared action and computes a scope hash over the sender,
  recipient, destination mother, title and personal space.
  \item The private phone displays that exact scope.  Sending requires a new signed approval whose
  scope hash still matches the prepared task.
  \item The sender mother encrypts and signs one task packet and records a private delivery receipt.
  \item The recipient mother verifies and decrypts the packet, then imports one pending delegation
  into the recipient's established Pilot task store.
  \item The recipient independently accepts or declines from a phone holding
  \texttt{task.respond}.  Acceptance makes the local task open; decline leaves it declined.
  \item The recipient mother signs a content-free response receipt.  The sender mother verifies
  the peer signature and reconciles its copy to the same state.
\end{enumerate}

Knowledge of the sender's contact name is not sufficient authority, and receipt of a task is not
treated as acceptance.  A shared screen may show a privacy-safe count but cannot expose or accept
the private request without the recipient's active identity and phone authority.

\subsection{Replay, isolation and failure behaviour}

Durable idempotency records prevent duplicate sender actions across retries.  Durable packet
identifiers prevent duplicate imports after reconnection.  A repeated valid packet returns the
existing object instead of creating a second message or task.  Reusing an idempotency key for
different content, modifying an encrypted packet, using an untrusted mother, reading another
persona's stream or acting from a revoked phone fails closed.

The automated vertical-slice exercise creates two isolated mothers, two private identities and two
independently paired phones.  It sends an encrypted task from Alice to Bob, proves that neither the
sender's persisted transport state nor the network packet contains the task title, verifies the
single import, accepts it on Bob's phone and reconciles the signed receipt back to Alice.  Separate
tests cover encrypted typed messages, duplicate retries, transport tampering, cross-person access
and lost-phone revocation.

\subsection{Remaining Stage 4 work}

The coded vertical slice is operational, but Stage~4 is not yet formally closed.  One field gate
remains explicitly open:
\begin{itemize}
  \item exercise text, voice-note and push-to-talk reconnection without duplication on real phones.
\end{itemize}

PHO, PhotonPay, wallets, balances, staking, bonds and minting remain excluded from the Unified
Inbox feature policy and routes.  The encrypted communications foundation does not activate the
dormant economic layer.

\subsection{Live mobile binding}

The mobile reference application is now bound to the production pairing boundary rather than a
visual pairing simulation.  The phone first retrieves the mother's short-lived discovery
descriptor over trusted HTTPS and verifies its Ed25519 signature, endpoint and fingerprint.  It
then creates or reuses a non-extractable browser-held Ed25519 possession key and requests one
five-minute challenge for the active persona.  The mother displays the six-digit confirmation
locally; the API never returns that code to the phone.

After the user checks the displayed trust words and enters the local code, the phone signs the
complete challenge.  The mother verifies possession, consumes the challenge exactly once and
returns a signed phone certificate plus an eight-hour, least-privilege lease.  The phone verifies
both mother signatures before retaining the connection in protected IndexedDB storage.  A copied
code without the original phone key cannot complete the exchange.

The same HTTPS service now exposes a persona-bound Unified Inbox stream.  A request must contain
the signed certificate and current lease and must carry \texttt{inbox.read}; the mother rejects
expired, revoked, cross-person or scope-incomplete connections.  Only after this check does the
mobile Inbox replace preview cards with that person's decrypted messages, task proposals and
content-free receipts.  If discovery, certificate validation, lease validation or Inbox loading
fails, the interface says that the preview remains active or that reconnection is required.  It
does not present fixture content as a live personal stream.

This completes the secure polling boundary and the phone-to-Inbox product binding.  It does not
close Stage~4: encrypted attachments, unified search, voice notes and push-to-talk remain
intentionally open.

\subsection{Authenticated live delivery and reconnection}

The Inbox now has a dedicated TLS WebSocket channel alongside the signed HTTPS action boundary.
The discovery descriptor advertises the live address only when the service is active, and the
phone accepts it only when it uses \texttt{wss}, names the same verified mother host and exposes
the exact versioned Inbox path.

Opening the socket is not authentication.  The first frame contains a narrowly defined request
with the persona, device, lease, fresh nonce and last observed Inbox revision.  The phone signs
that request with its non-extractable possession key.  The mother validates the certificate,
active lease, \texttt{inbox.read} scope, identity binding, device binding, lease binding and phone
signature before sending any private item.  An unknown origin, altered request, copied
certificate without its key, expired lease or revoked phone closes with a policy violation.

After authentication the mother sends only revisions newer than the phone's cursor.  The phone
retains the last applied revision during a connection, reconnects with bounded exponential
backoff and supplies that cursor in its newly signed opening request.  The server reports the
resume revision and explicitly records zero replayed duplicates.  A separate signed polling
request remains available as a recovery path when WebSocket transport is unavailable.

Live sockets are read-only notification surfaces.  They cannot send a message, accept a task or
perform an external action.  Every mutation continues through a separately scoped, signed and
idempotent action request.  This separation prevents a long-lived display connection from
becoming ambient authority.

\subsection{Bilateral task lifecycle}

The structured-task flow now continues beyond first delivery and acceptance.  A sender may
correct the task title or cancel the request while it remains active.  Those changes are carried
to the recipient mother as new authenticated and encrypted task-update packets.  The recipient
verifies the sending mother, destination persona and original action relationship before changing
either the private Unified Inbox projection or the authoritative Pilot task record.

The recipient may mark the request read, accept or decline it, snooze an accepted task, complete
it, report a bounded failure or record expiry after the declared deadline.  These state changes
return to the sender as Ed25519-signed receipts.  Receipts contain the action identifier, state,
time and evidence hash but never the task title or other private content.  The sender verifies the
trusted peer signature and destination mother before reconciling its own copy.

All phone mutations use a canonical signed request and a durable idempotency key.  Repeating the
same request returns its original receipt; reusing its key for different content fails closed.
Packet identifiers similarly suppress repeated imports after transport reconnection.  Automated
two-mother tests cover correction, duplicate update delivery, read, accept, snooze, completion,
cancellation, failure, expiry, altered receipts and synchronization with the existing personal
task authority.  The reference phone interface exposes only state-appropriate controls so that a
recipient cannot correct the sender's request and a sender cannot accept work on the recipient's
behalf.

\subsection{Encrypted attachments and unified private search}

The same Inbox can now carry bounded document, text and image attachments.  Before sending, Pilot
normalizes the filename, accepts only an explicit media-type allowlist, rejects empty or oversized
content and computes a SHA-256 digest over the exact bytes.  The bytes are included only inside
the recipient-specific X25519 and AES-GCM packet.  The network envelope therefore exposes routing
metadata and ciphertext, not the file or its filename.

Each mother stores its local copy as a separate AES-256-GCM blob with owner-only filesystem
permissions.  The attachment identifier, owner and content hash form authenticated additional
data.  Opening a file requires a fresh request signed by a phone with \texttt{inbox.read}; Pilot
then confirms that the phone's persona is a participant in the parent message.  Authentication
failure, an altered blob, a changed network packet, a disallowed media type, an invalid filename
or a size above one megabyte fails closed.  The mobile client receives verified bytes only after
that check and creates a short-lived local download URL which it revokes after use.

Inbox search is similarly executed behind the persona-bound read authority.  It decrypts only
that person's bounded recent window and searches messages, task titles and attachment filenames
as one mixed result set.  Filters for all items, messages, tasks and files change the projection,
not the underlying data owner.  Results are capped at fifty and ordered by creation time.  No
plaintext search index, provider-side content copy or separate files Inbox is created.

Automated evidence covers encrypted storage and transport, successful recipient retrieval,
content-hash verification, local-blob and packet tampering, media-type rejection, cross-person
denial and mixed search filtering.  Actual peer delivery still uses the surrounding Pilot
transport adapter; adding a file does not grant the mobile display ambient send authority.

\subsection{Encrypted voice notes and production push-to-talk authority}

Voice notes now use the authenticated Inbox attachment channel rather than a parallel audio
store.  The sender supplies an allowlisted audio format, a duration between one tenth of a second
and two minutes, and no more than 1.5 megabytes of encoded audio.  Retention is constrained to a
minimum of five minutes and a maximum of seven days, with one day as the default.  The mother
encrypts the bytes before local persistence, transports them only inside the recipient-specific
authenticated packet, and refuses to release expired content.  On expiry its local encrypted
blob is removed; raw microphone audio is not written to logs or a public proof rail.

Push-to-talk is governed by a short-lived floor issued by the mother brain.  Acquisition names one
exact recipient and conversation, and the floor is bound to the active persona, signed phone
device identifier and capability-lease identifier.  A different browser, phone, persona or lease
cannot renew, release or attach audio to that floor.  The floor lasts fifteen seconds, can be
renewed only by its holder, expires before the next floor decision when abandoned, and is released
atomically when its voice note is accepted.  Durable idempotency prevents repeated press, renew,
release or send requests from creating duplicate state.

The mobile control makes this lifecycle visible.  Pressing and holding requests the governed
floor before opening the microphone; recording is visibly indicated and the floor renews during a
long hold.  Release stops the recorder and all microphone tracks, submits one encrypted note and
refreshes the Inbox.  Failure also stops local capture and relinquishes the floor.  Playback first
performs a fresh possession-authorized attachment read and uses a temporary local object URL.

\subsection{Structured reminder, follow-up, calendar and document-review cards}

Four additional private envelope types now share the same signed and encrypted mother-to-mother
path.  Reminder and follow-up cards require an exact due time.  Calendar cards require parseable
start and end times with the end strictly after the start.  Document-review cards carry a bounded
title and private detail without pretending that a document provider has been changed.  The
receiver verifies the encrypted card hash before rendering it in the unified chronological Inbox.

These cards are proposals, not execution receipts.  Their transport does not send an email,
create a provider calendar event, schedule an operating-system notification or modify a document.
Those effects remain subject to their respective provider adapters, exact private approval and
verified execution receipts.  Automated two-mother evidence covers encrypted voice delivery,
retention and duration bounds, floor replay, device and lease isolation, forced floor expiry and
replacement, all four card types, card search and invalid calendar scope.

Stage~4 therefore has all seventeen coded items implemented.  Formal closure still waits for one
real-phone field exercise in which text, voice-note and push-to-talk delivery reconnect without a
duplicate.  This unresolved physical test remains recorded rather than being replaced by a mock.
